% Be sure to set \appendix or \appendices in structural/body.tex
\chapter{Literature review}
\label{chap:Lit_Review}

As we mentioned before, this work focuses on a two-player zero-sum orbital Target-Attacker-Defender game in Low Earth Orbit (\acs{leo}); the research problem is more difficult than a standard single-agent guidance problem. This is because an agent's best action depends not only on the orbital dynamics and arena geometry, but also on the changing behavior of an intelligent adversary. As our framework accounts for both full-state and partial-state information cases, the learned policies must operate while interacting with a classical \acs{gnc} architecture, including the \acs{cbfqp} controller and, only in the partially observable setting, a bearings-only \acs{ekf} to estimate the opponent's state. With this problem setup in mind, the most relevant prior literature includes orbital pursuit-evasion games, differential games, reinforcement learning for adversarial planning, and estimation and control-constrained learning methods.

Orbital pursuit-evasion and adversarial guidance games have become a prevalent area of research in aerospace autonomy. Plenty of examples of related work have been published in recent years, mostly on optimal-control-based methods:

\cite{gong2020continuousthrust} developed a classical differential game of satellite pursuit-evasion under continuous thrust using a Clohessy--Wiltshire (\acs{cw}) model. Moreover, \cite{selvakumar2017orbitalplane} developed a two-player zero-sum differential game with input constraints, providing a semi-analytical capture-time solution to their problem using closed-form optimal controls. Then \cite{liu2025multipletoone} proposed a one-step look-ahead Nash equilibrium method for cooperative pursuit-evasion in cooperative spacecraft. And lastly, \cite{drucker2026lowthrustgames} framed a low-thrust spacecraft pursuit-evasion problem as a linear-quadratic zero-sum game with terminal speed constraints, from which they obtained guidance laws via cost-weight selection demonstrating improved performance against actively evading targets.


In addition to more classical approaches to these rigorous space-bound problems, there has been a rise in journal-published research trying to solve variations of these games but opting for online reinforcement learning (\acs{rl}) approaches in recent years. Despite the longstanding criticism that core reinforcement learning algorithms do not provide theoretical convergence guarantees for learned policies, researchers have found problems in which these methods have been useful and have provided empirically convergent policies, strong performance relative to their assigned tasks, and excellent computational efficiency.


\cite{li2025heterogeneousppo} successfully developed a \acs{ppo} framework for multi-agent orbital pursuit-evasion, outperforming shooting methods in their optimality index metrics and computational efficiency while providing robustness to suboptimal opponent strategies. Next, \cite{cao2026historyawareppo} proposed the \acs{hprrppo} algorithm to do alternating learning for non-fixed-duration turn-based orbital game problems, showing effective interception against evaders with compute speed suitable for real-time deployment and reported a compute runtime of 1 ms per decision. Furthermore \cite{zhang2022nearoptimal} used deep reinforcement learning to achieve near-optimal guidance laws in one-to-one orbital pursuit-evasion games. Finally, \cite{zhao2023prdmaddpg} proposed a \acs{maddpg}-based method for impulsive orbital pursuit-evasion.

The fact that all these papers above have been published in the past few years makes it clear
that there is growing interest in developing reinforcement learning-based methods for planning in adversarial orbital dynamics. This is especially relevant given the generalized growing interest within the autonomy community to implement machine learning-based frameworks to complement or even replace classical control architectures. 


Furthermore, additional work has been made in the area of estimation-aware \acs{rl} frameworks. While physical \acs{ai} is something roboticists and industry alike are pushing towards, one must remember that real-life robots will not have full-state measurements of their environment as we can achieve in a simulation. Instead, they need to use sensors to characterize their surroundings and plan accordingly. Generally speaking, this is problematic on the decision-making side, as it can turn a clean Markov Decision Process (\acs{mdp}) present in a full-state environment into a Partially Observable Markov Decision Process (\acs{pomdp}), which is harder to characterize (\cite{kaelbling1998pomdp}). 

\cite{ye2017despot} showed that modern decision-making under partial observability can be carried out by planning over an approximated belief tree, showing that near-optimal policies can be obtained in partial observability without solving for the full \acs{pomdp}. Then \cite{igl2018dvrl} showed reinforcement learning under partial observability using latent-state inference, meaning that a belief representation is learned from incomplete observations; their method outperforms baselines in partially observable tasks. Lastly, \cite{zhang2015continuousmemory} studied partially observable learning with continuous control, where they handled the problem by augmenting the policies with continuous memory states, finding that learned policies can memorize relevant information needed to act with incomplete observations.

Moreover, some works have tried to fill in the gaps between learning frameworks and classical control frameworks by combining the most desirable parts of each into a single integrated architecture.

\cite{berkenkamp2017safembrl} shows that learning can be constrained via classical stability/safety filtering, using Lyapunov safety certificates instead of leaving the learning unconstrained. \cite{koller2019safeexplorationrl} bridged safe exploration and reinforcement learning through model predictive control, showing that uncertainty-aware trajectories and terminal safety constraints can be made safely while improving learning performance. Most importantly for the sake of this work, \cite{emam2021saferlrcbf} showed how a Control-Barrier-Function (\acs{cbf}) safety layer can be used to modify the learned policy's behavior by projecting the policy actions in a way that induces a safe learning transition map, preserving the safety of the learned model while allowing a reward-based policy to be learned on the side.

We notice a few gaps in the existing literature. Namely, we see that most current orbital \acs{rl} work, including those works cited in this literature review, assume privileged, full-state information within their frameworks, and practically none of the presented works study learned policies with online estimation in the loop. Almost none of these papers study the interaction of \acs{rl} policies when intertwined with safety filters and classical control in the loop. 

The novelty we present in this work is the integration of these pieces into a coherent and stable \acs{parlpo} framework. For this, we use a combination of reinforcement learning techniques that can lead to coherent, robust learning strategies in zero-sum orbital defense game settings: using \acs{ppo} to find learned policies under the given set of hyperparameters and complex reward shaping, doing so without the help of an equilibrium/Nash prior solution, and testing how to implement this method under a partial-observability setting via bearings-only Kalman Filter measurements. In this work, we contribute to the present literature in the following ways:

\begin{itemize}
    \item Using a staged \acs{ppo} training pipeline instead of residual networks to train the agents' policies in a zero-sum, two-player orbital game. Namely, the defender and attacker policies are trained in alternating phases against frozen and randomized mixtures of opponent policies. The goal of this is to produce policies that become more specialized, yet more generalizable, to complete their given task within the game framework instead of overfitting to any one rival. 

    \item Extending the \acs{ppo} defense framework from the previous bullet to partial-observation training with bearings-only estimator measurements. In other words, this work evaluates policies where an agent does not have access to the states of their adversaries. Instead, these have to be measured via bearings-only Kalman-filter estimates, based on the work by \cite{aidala1979kalman}. This provides a controlled test on how learned orbital defense policies degrade under sensor and partial observability limitations that would be present in a real-world implementation.
    \item Studying the interaction between learned policies in full and partial observability within the context of a classical \acs{cbfqp} safety filter. This shows how a model-based classical controller modifies or constrains learning behavior during training and rollouts.
    \item Evaluating if policies trained under \acs{hcw} linear time-invariant relative orbital dynamics can transfer to elliptical dynamics that were not seen during training. This should provide a controlled test to see if policies retain coherent behavior and performance when deployed outside the exact orbital dynamics model used during training.
\end{itemize}

So in summary, existing literature shows that orbital pursuit-evasion remains a relevant \acs{gnc} problem to research and that reinforcement learning is increasingly being used by various works to obtain guidance policies. Nevertheless, most of the existing learning-based work cited above assumes access to privileged state information and separates the learned guidance problem from a spacecraft's estimation and safety-based controls architecture. The contributions made by this work are therefore the integration of all these algorithms into a phased learning framework.

Specifically, by developing the \acs{parlpo} framework, we research how alternating, \acs{ppo}-trained attacker and defender policies perform given opponent randomization, bearings-only estimator measurements, and \acs{cbfqp} safety control filtering. This positions the work done on \acs{parlpo} as an effective architecture and case study for work on learning-based orbital defense subject to more realistic constraints than the full-state baseline.
